% ===========================================================================
\section{完备性定理：亨金构造}\label{sec:henkin}

\begin{strand}
完备性定理断言 $\Gamma\models\varphi\Rightarrow\Gamma\vdash\varphi$：
凡被所有模型满足的，都可以形式推导出来。对 $\Gamma\cup\{\neg\varphi\}$ 取逆否，
定理化为一个更对称的形式——
\[
  \Gamma\nvdash\bot \quad\Longleftrightarrow\quad \Gamma\text{ 有模型},
\]
即\emph{一致的句子集必有模型}。从右到左是可靠性，对推导树作归纳即可；
本节处理从左到右。
\end{strand}

\block{困难所在}
一致性是纯语法性质，它只说不存在以 $\Gamma$ 为前提、以 $\bot$ 为结论的有限推导；
而一个模型需要两样语义对象：一组充当论域的元素，以及每个原子句子在其上的真值。
直接面对一般的 $\Gamma$，两者都没有着落。具体有两处缺口：
\begin{enumerate}[label=(\alph*),leftmargin=1.8em,itemsep=1pt]
  \item \emph{见证缺口}。若 $\exists x\,\varphi(x)$ 可由 $\Gamma$ 推出，模型中就应当有
        具体的元素满足 $\varphi$。亨金模型打算用闭项（其等价类）充当论域元素，
        但若语言中根本没有为此准备的常元，存在句子的真值就没有着落。
  \item \emph{真值缺口}。$\Gamma$ 一般不能判定每个句子：对许多 $\varphi$，
        既有 $\Gamma\nvdash\varphi$ 又有 $\Gamma\nvdash\neg\varphi$。
        以“是否属于 $\Gamma$”来裁定真值，对这样的 $\varphi$ 没有定义。
\end{enumerate}
亨金构造依次补上这两个缺口：先向语言中加入见证常元，再把理论扩充为极大一致集。

\begin{definition}[理论、亨金系统]
给定公理集 $A$，其\emph{理论}为 $T=\{\text{句子 }\varphi\mid A\vdash\varphi\}$。
称 $T$ 为\keyword{亨金的一阶系统}（Henkin theory），若对每个形为
$\exists x\,\varphi(x)$ 的句子，都存在常元符号 $c$（称为\emph{亨金见证}，
Henkin witness），使
\[
  \bigl(\exists x\,\varphi(x)\bigr)\longrightarrow\varphi(c)\ \in\ T.
\]
换言之，理论一旦断言某个性质有实例，语言中就备好了指称该实例的常元。
\end{definition}

\begin{theorem}[模型存在定理]\label{thm:completeness}
在上述含量词规则与等号规则的自然演绎系统中，一致的句子集必有模型。
由此即得完备性：$\Gamma\models\varphi$ 蕴含 $\Gamma\vdash\varphi$。
\end{theorem}

\begin{proof}
分三步。

\par\smallskip\noindent\textbf{第一步：亨金扩充（Henkin extension）}\quad
原理论一般不是亨金的，因此造递增链 $T=T_0\subseteq T_1\subseteq T_2\subseteq\cdots$：
$T_{i+1}$ 在 $T_i$ 的语言中对每个句子 $\exists x\,\varphi(x)$ 引入一个\emph{新}常元
$c_\varphi$，并加入公理
\[
  \exists x\,\varphi(x)\longrightarrow\varphi(c_\varphi).
\]
必须逐层迭代而不能一次完成：新常元进入语言后会产生新的存在句子，它们又需要
各自的见证。取 $T'=\bigcup_{i\ge0}T_i$，则 $T'$ 是亨金的。
一致性在扩充下保持：若某个 $T_i\vdash\bot$，推导只涉及有限多条新见证公理；
从新到旧逐条把这些公理还原为相应的存在句子（$\exists$ 消去与引入），
便在 $T$ 中复现了矛盾，与 $T$ 一致矛盾。

\par\smallskip\noindent\textbf{第二步：极大一致扩充}\quad
由 Lindenbaum 引理，$T'$ 包含于某个极大一致集 $T_{\max}$。该引理第一周已在
命题逻辑情形证明；一阶情形的论证相同：枚举全部句子，逐个吸收且保持一致，
推导的有限性保证并集仍一致。极大一致性给出两条关键性质：每个句子 $\varphi$
与 $\neg\varphi$ 恰有一个属于 $T_{\max}$；且 $T_{\max}$ 对推导封闭。
真值缺口由此补上：每个句子的“真值”现在都有定义——它是否在 $T_{\max}$ 中。

\par\smallskip\noindent\textbf{第三步：构造模型}\quad
论域取亨金常元的\keyword{等价类}：
\[
  \Omega=\bigl\{[c_\varphi]\ \big|\ \exists x\,\varphi(x)\in T_{\max}\bigr\},
  \qquad
  c_\varphi\sim c_\psi \ \Longleftrightarrow\ (c_\varphi=c_\psi)\in T_{\max}.
\]
由 RE1--RE3，$\sim$ 是等价关系。函数与谓词按语法方式解释：
$f^{\M}([c_1],\ldots,[c_k])=[f(c_1,\ldots,c_k)]$，以及
$([c_1],\ldots,[c_k])\in P^{\M}$ 当且仅当 $P(c_1,\ldots,c_k)\in T_{\max}$；
解释的良定义性（与代表元选取无关）由代换规则 $\mathrm{RE}_4$ 保证。
至此每个原子句子的真值都已确定。对公式作结构归纳，证明\emph{真值引理}
\[
  \M\models\varphi \quad\Longleftrightarrow\quad \varphi\in T_{\max}.
\]
原子情形由定义；联结词情形由 $T_{\max}$ 的封闭性与完备性（$\varphi$ 与
$\neg\varphi$ 恰居其一）。关键是存在量词情形：若 $\exists x\,\varphi(x)\in T_{\max}$，
见证公理给出 $\varphi(c_\varphi)\in T_{\max}$，由归纳假设得
$\M\models\varphi(c_\varphi)$，故 $\M\models\exists x\,\varphi(x)$；
反之若 $\exists x\,\varphi(x)\notin T_{\max}$，则 $\neg\exists x\,\varphi(x)\in T_{\max}$，
由归纳假设与封闭性，不存在满足 $\varphi$ 的论域元素，故
$\M\not\models\exists x\,\varphi(x)$。这正是第一步加入见证常元的原因：
没有它们，$\M$ 中可能没有元素来实现理论断言的存在性。

于是 $\M\models T_{\max}\supseteq\Gamma$，即一致的 $\Gamma$ 有模型。
\end{proof}

\begin{corollary}[紧致性定理]
若句子集 $\Gamma$ 的每个有限子集都有模型，则 $\Gamma$ 有模型。
\end{corollary}
\begin{proof}
形式推导只含有限条前提，故“每个有限子集都有模型”经可靠性给出
“每个有限子集都一致”，进而 $\Gamma$ 一致；再由定理~\ref{thm:completeness} 得证。
\end{proof}

\begin{remark}\label{rem:countable}
亨金模型的论域由语言中的常元之等价类构成。语言至多可数时，常元至多可数，
故模型至多可数。下一节将用到这一观察。
\end{remark}
